\section{Développement}\label{sec:dev}
\subsection{Réalisation des versions}
Les cinq versions, les scripts de test et le programme d'analyse ont été produits avec l'assistance d'un outil
d'intelligence artificielle (Claude, d'Anthropic), ce que l'énoncé du travail autorise. Chaque version reprend la
même structure : un routeur, la validation des entrées, des requêtes SQL identiques et la construction de la réponse
JSON. Toutes utilisent le mode WAL de SQLite \cite{sqlitewal}, quatre connexions à la base (une seule pour Node.js, dont le code
s'exécute dans un seul fil) et des connexions HTTP persistantes.
\begin{itemize}[leftmargin=*,itemsep=0pt,topsep=2pt]
\item \textbf{Python} : Flask, avec le serveur de développement de Werkzeug en mode multifil (HTTP/1.1) et le module
\texttt{sqlite3} ;
\item \textbf{Node.js} : module \texttt{http} natif et \texttt{node:sqlite}, requêtes préparées une fois ;
\item \textbf{PHP} : serveur intégré \texttt{php -S} avec quatre processus de travail, PDO SQLite à connexion persistante ;
\item \textbf{Go} : \texttt{net/http} et un accès SQLite écrit avec cgo (liaison à \texttt{libsqlite3}) ;
\item \textbf{Rust} : bibliothèque standard seulement : serveur HTTP/1.1 minimal à un fil par connexion, lecteur
JSON écrit à la main et appels FFI à \texttt{libsqlite3}.
\end{itemize}
La machine n'ayant pas d'accès au réseau, les frameworks cités dans l'énoncé à titre d'exemple (Axum, Express,
Laravel) n'ont pas pu être installés. Cela signifie que les versions Go et Rust utilisent un serveur plus
rudimentaire que celui d'un projet réel, et que Python (Flask) est la seule version bâtie sur un framework.

\subsection{Équivalence fonctionnelle}
L'équivalence est vérifiée par un scénario de 69 requêtes rejoué sur chaque version, à partir d'une base neuve : 26
réponses 200, 3 créations (dont un titre avec accents, guillemets et caractères japonais), 1 suppression, 26 erreurs
400 (JSON invalide, statut ou priorité hors domaine, page ou limite non numérique, clé étrangère inexistante,
titre de 201 caractères) et 13 erreurs 404. Pour chaque requête, le code HTTP et le corps JSON analysé sont comparés.
Les cinq versions donnent exactement les mêmes réponses sur les 69 requêtes. Pendant les essais de charge, les
codes de réponse sont aussi contrôlés : le nombre total d'erreurs est indiqué dans la section~\ref{sec:res}.

\subsection{Difficultés rencontrées}
\emph{Versions de SQLite.} Python et Node.js embarquent leur propre SQLite (3.53), alors que PHP, Go et Rust utilisent
la bibliothèque du système (3.45.1). Cet écart n'a pas pu être supprimé ; il est une limite de l'étude (section IV-C).
\emph{Représentation du JSON en Go.} Une première version Go construisait des \texttt{map} génériques ; son débit
sur F2 était plus faible. Remplacer ces \texttt{map} par des structures typées a rendu l'implémentation plus
idiomatique, et l'équivalence a été revérifiée.
\emph{Serveurs « de développement ».} Le serveur Werkzeug de Flask et \texttt{php -S} sont déconseillés en production ;
les résultats de Python et PHP concernent donc cette configuration, non un déploiement avec Gunicorn ou
PHP-FPM.
\emph{Connexions persistantes.} Flask répond en HTTP/1.0 par défaut ; nous avons forcé HTTP/1.1 pour que toutes les
versions réutilisent leurs connexions.
